\section{Przegląd literatury}
\label{sec:literatura}

W niniejszej sekcji przedstawiamy genezę algorytmu Monte Carlo Tree Search,
najważniejsze rozszerzenia istotne dla tego projektu oraz literaturę
dotyczącą gry Breakthrough w kontekście MCTS.

\subsection{Geneza MCTS i UCT}

Metoda Monte Carlo Tree Search została zaproponowana przez Coulom\,a
\cite{coulom2006efficient} jako algorytm przeszukiwania dla gry Go,
gdzie wcześniejsze podejścia oparte na alpha-beta nie potrafiły poradzić
sobie z dużym branching factor (typowo 200--400 ruchów w pozycji).
Kluczowa innowacja Coulom\,a: zamiast przeszukiwać drzewo gry do pełnej
głębokości i ewaluować liście funkcją heurystyczną, rozwijano drzewo
selektywnie i zamiast ewaluacji symulowano losowe partie do końca,
agregując wyniki w statystykach Monte Carlo per stan.

W tym samym roku Kocsis i Szepesvári~\cite{kocsis2006bandit} wprowadzili
algorytm UCT (Upper Confidence Bounds applied to Trees), łączący metodę
Coulom\,a z regułą wyboru z teorii bandytów wieloramiennych. Funkcja
selekcji UCB1 daje matematycznie ugruntowany kompromis między
eksploatacją węzłów już dobrze zbadanych a eksploracją mniej odwiedzonych.

Kompleksowy przegląd MCTS i jego wariantów do roku 2012 przedstawili
Browne et al.~\cite{browne2012survey} ze szczególnym uwzględnieniem
zastosowań poza Go: gry kart, gry planszowe, planowanie. Praca ta jest
podstawowym punktem odniesienia metodologicznego dla niniejszego projektu.

\subsection{Rozszerzenia MCTS istotne dla projektu}

\paragraph{RAVE (Rapid Action Value Estimation).}
Gelly i Silver~\cite{gelly2011monte} wprowadzili technikę RAVE oraz
powiązane statystyki AMAF (All-Moves-As-First), będące rozszerzeniem
wcześniejszych prac nad polityką uczenia w grze Go. Główna idea: jeśli
ruch $m$ okazywał się dobry w rolloutach wychodzących z poddrzewa węzła
$n$, to $m$ jest również prawdopodobnie dobry z samego $n$, nawet jeśli
$m$ nie był jeszcze próbowany jako pierwszy. Selekcja RAVE blenduje
oszacowania UCT i AMAF z wagą $\beta$ malejącą wraz z liczbą odwiedzin
węzła --- oryginalna formuła $\beta(n, \tilde{n})$ uwzględnia zarówno
liczbę odwiedzin węzła $n$, jak i liczbę odwiedzin AMAF $\tilde{n}$.

W~grze Go RAVE daje silne przyspieszenie zbieżności --- Gelly i Silver
raportują 20--40\% wzrost win-rate przy tym samym budżecie iteracji.
Kluczowe założenie sukcesu: \emph{translation invariance} wartości ruchu ---
postawienie kamienia w punkcie $(x, y)$ ma podobną wartość w różnych
kontekstach taktycznych.

\paragraph{Progressive Bias / Progressive Strategies.}
Chaslot et al.~\cite{chaslot2008progressive} zaproponowali rodzinę
,,progressive strategies'' wstrzykujących wiedzę domenową do MCTS.
Centralna z technik --- Progressive Bias --- dodaje do funkcji selekcji
człon $H(m) / (n + 1)$, gdzie $H(m)$ jest oszacowaniem heurystycznym
ruchu $m$. Człon ten dominuje wczesne iteracje (gdy $n$ jest małe),
po czym zanika z rosnącym $n$, oddając kontrolę statystykom UCB1.
Chaslot pokazał, że taka kombinacja przewyższa zarówno czyste UCT,
jak i czystą heurystykę, w grze Go i innych domenach.

\paragraph{MAST (Move-Average Sampling Technique).}
Finnsson i Björnsson~\cite{finnsson2008simulation} w kontekście General Game
Playing zaproponowali MAST --- globalną tabelę średnich wyników per ruch,
aktualizowaną po każdej symulacji. Ta tabela następnie wpływa na rolloutsy
(softmax/$\epsilon$-greedy zamiast uniform random). MAST był rozważany
jako alternatywna druga modyfikacja UCT (zamiast Progressive Bias),
ostatecznie odrzucony jako mniej spójny tematycznie z heurystyką H3.

\paragraph{Pozostałe rozszerzenia.}
Lorentz~\cite{lorentz2016early} opisuje \emph{Early Playout Termination}:
wczesne kończenie symulacji w stanach o oczywistym zwycięzcy, znacząco
przyspieszające MCTS w grach z wyraźną strukturą terminalu (jak Breakthrough).
Inne rozszerzenia warte uwagi (poza zakresem tego projektu): UCT z transposition
tables, virtual loss dla zrównoleglenia, decisive/anti-decisive moves
(Teytaud \& Teytaud, AAAI 2010).

\subsection{Breakthrough w literaturze MCTS}

\paragraph{Status rozwiązaniowy gry.}
Saffidine et al.~\cite{saffidine2011solving} opracowali techniki ,,race
patterns'' i Job-Level Proof Number Search, pozwalające im rozwiązać
warianty Breakthrough do planszy 6$\times$5 włącznie; we wszystkich
\emph{wygrywa pierwszy gracz}. Niedawno Björnsson et al.~\cite{bjornsson2026solving}
rozwiązali również planszę 6$\times$6 (także na korzyść pierwszego gracza).
Klasyczna plansza 8$\times$8 pozostaje natomiast nierozwiązana --- nie ma
znanego solwera o ograniczonych zasobach pamięciowych zdolnego do pełnego
przeszukania jej drzewa. To czyni Breakthrough 8$\times$8 atrakcyjnym
benchmarkiem: nie istnieje wyrocznia ,,perfekcyjnego gracza'' definiująca
sufit jakości, więc porównanie wariantów MCTS ma realną wartość badawczą.

\paragraph{Heurystyki ewaluacji.}
Lorentz~\cite{lorentz2011breakthrough} (w pracy nad EinStein würfelt nicht!,
ale stosowanej również do Breakthrough) opisuje proste heurystyki ewaluacji
pozycji zawierające zaawansowanie pionków, materiał oraz bezpieczeństwo
struktury (pionki na krawędziach są bezpieczniejsze, ponieważ mogą zostać
zbite tylko z jednej strony zamiast dwóch). Te elementy stanowią podstawę
naszej implementacji \texttt{HeuristicAgent}.

\paragraph{Branching factor i głębokość.}
W typowej pozycji Breakthrough 8$\times$8 liczba legalnych ruchów wynosi
około 15--30, a głębokość partii (przy losowej grze) to 40--80 półruchów.
Z naszych eksperymentów (sekcja~\ref{sec:wyniki}) średnia długość partii
UCT vs UCT na planszy 8$\times$8 wynosi ok.\ 57 ruchów. Branching factor
zalicza Breakthrough do gier ,,małych'' z punktu widzenia MCTS ---
znacznie mniej niż Go (200--400), porównywalnie do Hex (odpowiadający
rozmiar $11 \times 11$ ma branching ok. 80--100).

\paragraph{Status badawczy Breakthrough 8$\times$8.}
Klasyczna plansza 8$\times$8 pozostaje \emph{nierozwiązana} --- mimo postępów
w rozwiązywaniu mniejszych wariantów (do 6$\times$5 \cite{saffidine2011solving}
oraz 6$\times$6 \cite{bjornsson2026solving}), dla 8$\times$8 nie istnieje znany
,,perfekcyjny gracz''. Brak odniesienia do optimum czyni z 8$\times$8
wdzięczny poligon dla \emph{porównawczych} badań empirycznych nad wariantami
MCTS, gdzie miarą jest względna siła agentów, a nie odległość od gry idealnej.
